%% =====================================================================
%%  B9-T-18-05 -- what B9 contributes to "The Man in the Gray Coat"
%%  -------------------------------------------------------------------
%%  LAYER A: APPEND-ONLY. Written once at the entry's write, then left
%%  alone. If the source has more to say later, add a bullet -- do not
%%  rewrite. Material found ONLY here goes in the `unique` slot with
%%  the unique flag set, which makes it undeletable (rule R10).
%% =====================================================================
%% @kind: source
%% @id: B9-T-18-05
%% @book: B9
%% @topic: T-18-05
%% @locator: JASP 67(4):371--378; Method, Results, Discussion (13-feature analysis)
%% @status: distilled
%% @unique: yes
%% @integrated: yes
%% @updated: 2026-09-20

\dossier{B9}{T-18-05}{\bookshort{B9}}{JASP 67(4):371--378; Method, Results, Discussion}

\begin{distilled}
  \item The sample: 40 males aged 20--50, New Haven area, recruited by newspaper advertisement, paid \$4.50 and told the money was theirs no matter what; roles of teacher and learner assigned by a drawing rigged so the subject was always the teacher.
  \item The instrument: a simulated shock generator, 30 switches, 15-volt steps from 15 to 450, labelled Slight Shock through Danger: Severe Shock to XXX; each subject first took a 45-volt sample shock on his own wrist.
  \item Authority was carried entirely by symbols and manner: Yale's interaction laboratory, a 31-year-old biology teacher in a gray technician's coat, impassive, reading scripted prods; the learner was a mild-mannered 47-year-old accountant.
  \item The prods, used in sequence and renewed at each refusal: (1) Please continue; (2) The experiment requires that you continue; (3) It is absolutely essential that you continue; (4) You have no other choice, you must go on. Special prods: shocks are painful but cause no permanent tissue damage; whether the learner likes it or not, you must go on.
  \item Protest schedule: the learner pounds the wall at 300 volts, is silent from 315; the experimenter directs silence to be scored as a wrong answer.
  \item The result: 26 of 40 (65 percent) obeyed to the 450-volt maximum. The 14 defiers broke at 300 (5), 315 (4), 330 (2), 345, 360, 375 (1 each). No subject stopped before the victim's first protest at 300.
  \item The prediction gap: fourteen Yale psychology majors, given full details, estimated that 0--3 percent of 100 hypothetical subjects would go to the maximum (mean 1.2 percent); one-way-mirror observers also systematically underestimated obedience.
  \item The cost to the obedient: sweating, trembling, stuttering; nervous laughter in 14 of 40; full seizures in 3 --- one, a 46-year-old encyclopedia salesman, violent enough to halt the experiment. Subjects' post-run pain rating of the final shocks: mean 13.42 on a 14-point scale (they believed it was real).
  \item Payment was not the engine: 43 unpaid Yale undergraduates run separately obeyed at very similar rates (footnote 4).
  \item Milgram's 13-feature situational analysis: Yale sponsorship; a worthy purpose (knowledge); the victim volunteered; the subject volunteered; payment; a fair lottery for roles; ambiguous psychologist prerogatives in a closed setting; the ``painful but not dangerous'' assurance; the victim still playing the game until 300 volts; a conflict resolvable only by visible action; the asymmetry of abstract science against felt suffering; little time for reflection under persistent conflict; and two dispositions in collision --- not harming others versus obeying legitimate authority.
\end{distilled}
\begin{unique}
  \item The obedience distribution (65 percent full obedience; every refusal after the victim's 300-volt protest), the 1.2 percent prediction gap, the four-prod script with its exact wording, and the 13-feature situational analysis are unique to B9 in this corpus --- no other source measures obedience to a destructive command.
\end{unique}
